\chapter{Compactness}\label{ch:compact}

\begin{center}\itshape
On a closed interval a continuous function reaches a highest value. On an open interval it can keep climbing and never arrive. A bounded sequence of real numbers has a subsequence that settles. A bounded sequence of continuous functions can fail to settle in the uniform distance. The good cases share one property. How far does that property go?
\end{center}

Picture a walker who paces a closed stretch of road for as long as you like. Some spot has to see the walker come back, again and again, as close as you ask. On a closed segment that is the Bolzano--Weierstrass theorem, and it is why Volume~1 could locate maxima of continuous functions. On an endless road the walker can leave every neighbourhood and stay away. The closed unit ball in a space of functions has infinitely many directions, and the walker can leave in a new direction each time. Compactness is the name of the property that keeps the picture finite.

Finite sets are easy to handle. A function on a finite set has a largest value, and a finite set of positive numbers has a positive lower bound. Compactness lets an infinite space behave like a finite one in those two arguments. The definition took a long time to settle. Heine (1872) and Borel (1895) saw that a cover of $[a,b]$ by open intervals can be cut down to a finite subcover. Fr\'echet coined the word ``compact'' in 1906. Alexandroff and Urysohn (1923) took the cover property as the definition. The definition looks odd at first sight. It is exactly the move the proofs use. The same wording applies when a product has arbitrarily many factors.

The first job of this chapter is to keep the last promises of Volume~1: extreme values, Heine's uniform continuity, and Rolle's theorem with the mean value theorem. After that come the three compactness theorems analysis uses most often. Tychonoff's theorem says that a product of compact spaces is compact, with no restriction on the number of factors. In 1950 John Kelley showed that this statement is equivalent to the axiom of choice. Arzel\`a--Ascoli describes the compact sets of continuous functions. Stone--Weierstrass says when a family of functions can approximate every continuous function.

\section{Compact spaces}

\begin{definition}
An \emph{open cover} of a topological space $X$ is any family of open sets with union equal to $X$. The space $X$ is \emph{compact} when every open cover has a finite subfamily that still covers $X$. That finite subfamily is a \emph{finite subcover}. A subset $K\subseteq X$ is compact when it is compact in the subspace topology. Equivalently, every cover of $K$ by open sets of $X$ has a finite subcover.
\end{definition}

The same fact, read on complements, is often easier to use. The space $X$ is compact if and only if every family of closed sets with empty total intersection has a finite subfamily with empty intersection. A family of sets has the \emph{finite intersection property} when every finite subfamily has nonempty intersection. Thus $X$ is compact if and only if every family of closed sets with the finite intersection property has nonempty intersection. One nested case comes up often. If $F_1\supseteq F_2\supseteq\cdots$ are nonempty closed sets in a compact space, the whole sequence has a point in common.

\begin{think}
Test three spaces. Is $[0,1[$ compact? Is $\Z$, with the discrete topology, compact? Is every finite space compact?
\end{think}

The sets $[0,1-\frac1n[$ cover $[0,1[$, and no finite subcollection covers it: the rightmost of them stops short of $1$. The singletons cover $\Z$, and each point needs its own set. Every finite space is compact, and so is any space that has only finitely many open sets.

\begin{figure}[ht]
\centering
\fig{ch04-rw-subcover}
\caption{Three open intervals already cover $[0,1]$. The dashed interval belongs to the cover and can be dropped. Compactness of the segment is the statement that every open cover thins out in this way.}
\label{fig:subcover}
\end{figure}

Figure~\ref{fig:subcover} shows a single cover. The intervals $U_1$, $U_2$ and $U_3$ overlap, and together they contain $[0,1]$. The spare interval $V$ sits inside $U_2$. Heine--Borel, proved below, says that every open cover of the segment has some finite subcover of this kind.

\begin{proposition}\label{prop:compactbasic}
\begin{enumerate}[label=(\alph*)]
\item Every closed subset of a compact space is compact.
\item In a Hausdorff space, every compact subset is closed.
\item The continuous image of a compact space is compact.
\item Whenever the domain is compact and the target is Hausdorff, a continuous bijection is a homeomorphism.
\item A compact Hausdorff space is normal. Any two disjoint closed sets have disjoint open neighbourhoods.
\end{enumerate}
\end{proposition}
\begin{proof}
(a) Let $F\subseteq X$ be closed, and let $\mathcal U$ be a cover of $F$ by open sets. Then $\mathcal U\cup\{X\setminus F\}$ covers $X$. A finite subcover of $X$, with $X\setminus F$ removed if it was used, still covers $F$.

(b) Let $K$ be compact in a Hausdorff space $X$, and take $x\notin K$. For each $y\in K$, Hausdorff separation gives disjoint open sets $U_y\ni x$ and $V_y\ni y$. Finitely many of the $V_y$, say $V_{y_1},\dots,V_{y_m}$, cover $K$. The intersection $U=U_{y_1}\cap\cdots\cap U_{y_m}$ is an open neighbourhood of $x$. It misses every $V_{y_i}$, hence misses $K$. So $X\setminus K$ is open.

(c) Let $f:X\to Y$ be continuous, and let $\mathcal V$ be a cover of $f(X)$ by open sets. The preimages $f^{-1}(V)$, for $V\in\mathcal V$, cover $X$. Finitely many of them cover $X$, and the corresponding sets $V$ cover $f(X)$.

(d) Let $f:X\to Y$ be a continuous bijection onto a Hausdorff space. If $F$ is closed in $X$, then $f(F)$ is compact by (a) and (c), and closed in $Y$ by (b). Since $(f^{-1})^{-1}(F)=f(F)$, the inverse is continuous.

(e) First let $A$ be a single point $a$, and let $B$ be closed with $a\notin B$. By (a), the set $B$ is compact. The argument of (b) produces an open set $U\ni a$ and an open set $V=V_{y_1}\cup\cdots\cup V_{y_m}\supseteq B$ with $U\cap V=\varnothing$. Now let $A$ and $B$ both be closed and disjoint. For each $a\in A$, the previous case gives disjoint open sets $U_a\ni a$ and $V_a\supseteq B$. Finitely many of the $U_a$ cover the compact set $A$. Their union $U$ contains $A$, and the intersection $V$ of the matching sets $V_a$ contains $B$. Both are open, and they do not meet.
\end{proof}

Part~(d) is the test you will use most often. A continuous bijection from a compact space onto a Hausdorff space is a homeomorphism, and there is nothing further to check. Part~(d) makes the continuous bijection from the compact quotient $\R/\Z$ onto $S^1$ a homeomorphism. Chapter~\ref{ch:topology} also checks the inverse by hand. The usual maps from $[0,1]$ onto $\R/\Z$ and onto $S^1$ are continuous and onto, and they send two endpoints to one point, so they are not bijections. Part~(e), together with the remark after Proposition~\ref{prop:urysohn}, produces Urysohn functions on a compact Hausdorff space.

\begin{proposition}\label{prop:finiteproduct}
If $X$ and $Y$ are compact, then $X\times Y$ is compact. A finite product of compact spaces is therefore compact.
\end{proposition}
\begin{proof}
It is enough to treat covers by basic open sets $U\times V$. Given an arbitrary open cover $\mathcal W$, the basic open sets that sit inside some member of $\mathcal W$ still cover $X\times Y$. If finitely many of those basic sets cover the product, the finitely many members of $\mathcal W$ that contain them cover it as well.

Let $\mathcal B$ be a cover of $X\times Y$ by basic open sets, and fix $x\in X$. The slice $\{x\}\times Y$ is homeomorphic to $Y$, so it is compact. Finitely many members of $\mathcal B$ cover the slice. Drop any that misses the slice, and relabel the rest $U_1\times V_1,\dots,U_m\times V_m$. Each of them contains a point $(x,y)$, so $x\in U_j$. The intersection $U_x=U_1\cap\cdots\cap U_m$ is an open neighbourhood of $x$. The whole tube $U_x\times Y$ is covered by the same $m$ sets. A point $(x',y)$ of the tube has $(x,y)\in U_j\times V_j$ for some $j$, and $x'\in U_x\subseteq U_j$, so $(x',y)\in U_j\times V_j$. Finitely many of the sets $U_x$ cover the compact space $X$. The corresponding finite list from $\mathcal B$ covers $X\times Y$.
\end{proof}

The payoff is the closed segment. A closed box in $\R^n$ is a product of segments, and the walker on a closed road has nowhere to escape.

\begin{theorem}[Heine--Borel]\label{thm:heineborel}
Every closed interval $[a,b]$ in $\R$ is compact. A subset of $\R^n$ is compact if and only if it is closed and bounded.
\end{theorem}
\begin{proof}
Let $\mathcal U$ be an open cover of $[a,b]$, and let $S$ be the set of points $x\in[a,b]$ for which $[a,x]$ has a finite subcover from $\mathcal U$. Then $a\in S$. Let $s=\sup S$, which exists by the least upper bound property, Theorem~\ref{thm:reals}. Some member $U\in\mathcal U$ contains $s$, and with it a relative interval $]s-\delta,s+\delta[\cap[a,b]$ for some $\delta>0$. Pick $x\in S$ with $x>s-\delta$. Adding $U$ to a finite subcover of $[a,x]$ covers $[a,\min(s+\frac\delta2,b)]$. That right-hand endpoint lies in $S$. If $s<b$, then $\min(s+\frac\delta2,b)>s$, which contradicts the choice of $s$ as a least upper bound. So $s=b$, and the same endpoint is $b$ itself. Thus $b\in S$, and $[a,b]$ has a finite subcover.

Now let $K\subseteq\R^n$ be closed and bounded, say $K\subseteq[-M,M]^n$. The topology of $\R^n$ is the product topology, by Proposition~\ref{prop:product}(b), so the cube $[-M,M]^n$ is compact by Proposition~\ref{prop:finiteproduct}. The set $K$ is closed in that cube, hence compact. Conversely, a compact set $K\subseteq\R^n$ is closed by Proposition~\ref{prop:compactbasic}(b), because $\R^n$ is Hausdorff. It is bounded because the open balls $B(0,k)$, $k\ge1$, cover it, and finitely many of them suffice.
\end{proof}

You can see the conclusion on a small function before using it in general. On $[0,1]$, the identity $x\mapsto x$ reaches $0$ and $1$. On $]0,1[$ the same formula reaches neither.

\begin{corollary}[extreme values]\label{cor:extreme}
A continuous real function on a nonempty compact space is bounded, and it attains its maximum and its minimum.
\end{corollary}
\begin{proof}
Let $f:K\to\R$ be continuous, with $K$ nonempty and compact. The image $f(K)$ is compact in $\R$ by Proposition~\ref{prop:compactbasic}(c), so Theorem~\ref{thm:heineborel} says it is closed and bounded. Choose points of $f(K)$ within $\frac1n$ of the supremum. That sequence converges to the supremum, and a closed set contains the limit, so the supremum belongs to $f(K)$. The same argument applies to the infimum.
\end{proof}

Volume~1 proved the extreme value theorem for a closed bounded interval. The argument above is the same fact for an arbitrary nonempty compact space. The first use promised in Chapter~\ref{ch:metric} is the comparison of norms.

\begin{corollary}\label{cor:norms}
Any two norms on $\R^n$, or on $\C^n$, are Lipschitz equivalent. Every finite-dimensional normed space is therefore complete, and its compact subsets are exactly its closed bounded subsets.
\end{corollary}
\begin{proof}
Let $N$ be a norm on $\R^n$, and let $(e_i)$ be the standard basis. Then
\[N(x)\le\sum_i|x_i|N(e_i)\le C\|x\|_\infty,\qquad C=\sum_i N(e_i).\]
The triangle inequality gives $|N(x)-N(y)|\le N(x-y)\le C\|x-y\|_\infty$, so $N$ is continuous for $\|\cdot\|_\infty$. The unit sphere $S=\{x:\|x\|_\infty=1\}$ is closed and bounded, hence compact. On $S$ one has $N>0$. Corollary~\ref{cor:extreme} supplies a constant $c>0$ with $N\ge c$ on $S$. Homogeneity extends the two inequalities to every $x$:
\[c\|x\|_\infty\le N(x)\le C\|x\|_\infty.\]
The complex case is the real case, read in $\R^{2n}$. A finite-dimensional normed space becomes $\R^n$ or $\C^n$ once a basis is chosen, and the given norm is carried to a norm $N$ on that coordinate space. Lipschitz equivalence does not change completeness, which is Proposition~\ref{prop:complete}(b), and it does not change which sets are compact, which is Theorem~\ref{thm:heineborel}.
\end{proof}

On $\R^2$ the comparison is visible by hand. Writing $\|x\|_1=|x_1|+|x_2|$ and $\|x\|_\infty=\max(|x_1|,|x_2|)$, you have $\|x\|_\infty\le\|x\|_1\le 2\|x\|_\infty$.

\section{Compact metric spaces}

In a metric space, compactness can be read in two other ways: through sequences, and through finite nets.

\begin{definition}
A metric space $X$ is \emph{sequentially compact} when every sequence has a subsequence that converges in $X$. It is \emph{totally bounded} when, for every $\varepsilon>0$, finitely many balls of radius $\varepsilon$ cover $X$. The centres of those balls form an \emph{$\varepsilon$-net}.
\end{definition}

\begin{figure}[ht]
\centering
\fig{ch04-net}
\caption{Total boundedness in a picture: at every scale $\varepsilon$, finitely many balls of that radius cover the set, and the dots at their centres are an $\varepsilon$-net. Compactness is this finite picture at every scale, together with completeness.}
\label{fig:net}
\end{figure}

Figure~\ref{fig:net} draws the $\varepsilon$-net from condition~(iii) of the theorem below. Finitely many centres, at every scale, together with completeness, replace the open-cover definition.

\begin{theorem}\label{thm:compactmetric}
For a metric space $X$, the following three statements are equivalent:
\begin{enumerate}[label=(\roman*)]
\item $X$ is compact;
\item $X$ is sequentially compact;
\item $X$ is complete and totally bounded.
\end{enumerate}
If $X$ is compact, every open cover $\mathcal U$ has a \emph{Lebesgue number}: some $\delta>0$ such that every ball of radius $\delta$ lies inside a single member of $\mathcal U$.
\end{theorem}
\begin{proof}
Proof of (i)$\Rightarrow$(ii). Let $(x_n)$ be a sequence in $X$. Call $x$ a \emph{cluster point} of the sequence when every ball about $x$ contains $x_n$ for infinitely many indices $n$. Suppose there is no cluster point. Then each $x$ has an open ball $B_x$ that contains $x_n$ for only finitely many $n$. Finitely many of the balls $B_x$ cover $X$, so only finitely many indices $n$ exist, which is absurd. Thus some cluster point $x$ exists. Choosing indices $n_1<n_2<\cdots$ with $d(x_{n_k},x)<\frac1k$ produces a subsequence converging to $x$.

(ii)$\Rightarrow$(iii). A Cauchy sequence that has a convergent subsequence converges to the same limit, so $X$ is complete. Suppose $X$ fails to be totally bounded. Then some $\varepsilon>0$ admits no finite cover by balls of radius $\varepsilon$. Choose $x_1$ arbitrarily, then $x_2$ outside $B(x_1,\varepsilon)$, then $x_3$ outside $B(x_1,\varepsilon)\cup B(x_2,\varepsilon)$, and so on. The resulting points satisfy $d(x_m,x_n)\ge\varepsilon$ whenever $m\ne n$, so no subsequence is Cauchy.

(iii)$\Rightarrow$(ii). Let $(x_n)$ be a sequence. Finitely many balls of radius $1$ cover $X$, so one of them, called $B_1$, contains $x_n$ for infinitely many $n$. Let $N_1$ be that infinite set of indices. Finitely many balls of radius $\frac12$ cover $X$, so one of them, called $B_2$, contains $x_n$ for infinitely many $n\in N_1$. Call that set of indices $N_2$. Continuing, you obtain infinite sets $N_1\supseteq N_2\supseteq\cdots$ with $x_n\in B_k$ whenever $n\in N_k$, and with the radius of $B_k$ equal to $\frac1k$. Choose $n_1<n_2<\cdots$ with $n_k\in N_k$. If $j\ge k$, then both $x_{n_j}$ and $x_{n_k}$ lie in $B_k$, so $d(x_{n_j},x_{n_k})<\frac2k$. The subsequence is Cauchy. Completeness makes it converge.

(ii)$\Rightarrow$(i). Assume sequential compactness, and first produce a Lebesgue number of an open cover $\mathcal U$. If none exists, then for each $n$ some ball $B(x_n,\frac1n)$ lies in no member of $\mathcal U$. A subsequence $x_{n_k}$ converges to a point $x$. That point lies in some $U\in\mathcal U$, and $B(x,r)\subseteq U$ for some $r>0$. For large $k$ one has both $d(x_{n_k},x)<\frac r2$ and $\frac1{n_k}<\frac r2$, so
\[B\Bigl(x_{n_k},\frac1{n_k}\Bigr)\subseteq B(x,r)\subseteq U,\]
which contradicts the choice of the ball. Thus some $\delta>0$ is a Lebesgue number of $\mathcal U$. Sequential compactness has already given total boundedness, so finitely many balls of radius $\delta$ cover $X$. Each of them lies in some member of $\mathcal U$, and those finitely many members cover $X$.
\end{proof}

A compact metric space is bounded. It is also \emph{separable}: the union, over $k\ge1$, of finite $\frac1k$-nets is a countable dense set. In $\R^n$, Theorem~\ref{thm:compactmetric} is the Bolzano--Weierstrass theorem. A bounded sequence has a convergent subsequence, because the closure of its range is closed and bounded, hence compact.

The same theorem names a fact that Chapter~\ref{ch:metric} proved by hand. Corollary~\ref{cor:Zpcompact} already proved that $\Z_p$ is sequentially compact. Theorem~\ref{thm:compactmetric} upgrades that to compactness: every open cover of $\Z_p$ has a finite subcover. The argument given there was the implication (iii)$\Rightarrow$(ii) in a special case. The finite nets were the $p^n$ balls of radius $p^{-n}$. Every closed ball of $\Q_p$ is compact for the same reason. Like $\R$, the field $\Q_p$ is \emph{locally compact}: every point has a compact neighbourhood.

\begin{figure}[ht]
\centering
\fig{ch04-rw-escape}
\caption{In $X=(0,1]$ the points $1,\frac12,\frac13,\ldots$ have no subsequence that converges inside $X$. Every subsequence runs toward $0$, and $0$ is not a point of $X$, so $X$ is not sequentially compact.}
\label{fig:rwescape}
\end{figure}

Figure~\ref{fig:rwescape} is the sequential test on a set that looks almost like a segment. Let $X=(0,1]$ with the usual distance, and let $x_n=\frac1n$. Any subsequence is of the form $\frac1{n_k}$ with $n_k\to\infty$, so it converges to $0$ in $\R$. A limit in $X$ would also be a limit in $\R$, and limits in $\R$ are unique, so the only candidate is $0$. That point is not in $X$. Theorem~\ref{thm:compactmetric} says $X$ is not compact. The same sequence is Cauchy and has no limit in $X$, so $X$ is not complete either. It is still totally bounded: a finite grid of spacing $\varepsilon$ covers it. Compactness fails because the limit is missing, not because the set is too large to net.

\begin{think}
The closed unit ball of $C([0,1])$, in the sup norm, is closed, bounded and complete. Is it compact?
\end{think}

It is not. Let $f_n$ be a tent of height $1$ with support $[\frac1{n+1},\frac1n]$. Then $\|f_n\|_\infty=1$. For $n\ne m$ the supports are disjoint, so $\|f_n-f_m\|_\infty=1$. No subsequence is Cauchy. The ball is complete, and it is not totally bounded. The tents are a special case of a theorem of F.~Riesz from 1918: the closed unit ball of a normed space is compact only when the space is finite-dimensional. Volume~8 proves that theorem. The tents above do not rely on it. Arzel\`a--Ascoli, in Section~\ref{sec:ascoli}, will say which sets of continuous functions are compact.

\begin{theorem}[Heine]\label{thm:heine}
A continuous map $f:X\to Y$ from a compact metric space into a metric space is uniformly continuous. For every $\varepsilon>0$ there is a $\delta>0$ such that $d(x,y)<\delta$ implies $d(f(x),f(y))<\varepsilon$.
\end{theorem}
\begin{proof}
Suppose not. Then some $\varepsilon>0$ has the following property: for every $n$ there are points $x_n,y_n$ with $d(x_n,y_n)<\frac1n$ and $d(f(x_n),f(y_n))\ge\varepsilon$. Sequential compactness gives a subsequence $x_{n_k}$ converging to some $x$, and then $y_{n_k}\to x$ as well. Continuity at $x$ forces both $f(x_{n_k})$ and $f(y_{n_k})$ to converge to $f(x)$, so $d(f(x_{n_k}),f(y_{n_k}))\to0$. That contradicts the gap of size $\varepsilon$.
\end{proof}

The square $x\mapsto x^2$ on $\R$, and the reciprocal $x\mapsto\frac1x$ on $]0,1]$, are continuous and fail to be uniformly continuous. Compactness of the domain is the hypothesis that cannot be dropped. Heine's theorem is what makes the integral of a continuous function exist, once a fine subdivision is chosen: on each piece, $f$ moves by less than $\varepsilon$. That construction is carried out in Chapter~\ref{ch:integral}. The estimate itself is the theorem just proved.

\section{Rolle and the mean value theorem}

Volume~1 proved the mean value theorem from extreme values. The derivative is the same one used throughout Volume~1. A function $f$ defined near $c$ is \emph{differentiable} at $c$ when $\frac{f(x)-f(c)}{x-c}$ has a limit $f'(c)$ as $x\to c$ with $x\ne c$.

\begin{theorem}\label{thm:mvt}
\begin{enumerate}[label=(\alph*)]
\item (Fermat) If $f:\,]a,b[\,\to\R$ has a maximum or a minimum at an interior point $c$, and if $f$ is differentiable at $c$, then $f'(c)=0$.
\item (Rolle) Suppose $f:[a,b]\to\R$ is continuous, differentiable on $]a,b[$, and $f(a)=f(b)$. Then $f'(c)=0$ for some $c\in\,]a,b[$.
\item (Mean value theorem) Suppose $f:[a,b]\to\R$ is continuous and differentiable on $]a,b[$. Then $f(b)-f(a)=f'(c)(b-a)$ for some $c\in\,]a,b[$.
\item On an interval $I$, the derivative controls the function. If $f'\ge0$ on the interior, then $f$ is non-decreasing. If $f'=0$ on the interior, then $f$ is constant. If $|f'|\le M$ on the interior, then $f$ is $M$-Lipschitz.
\end{enumerate}
\end{theorem}
\begin{proof}
(a) At a maximum, the difference quotient $\frac{f(x)-f(c)}{x-c}$ is $\le0$ when $x>c$, and $\ge0$ when $x<c$. The limit $f'(c)$ is therefore both $\le0$ and $\ge0$, so it is $0$. At a minimum, apply the same argument to $-f$.

(b) Corollary~\ref{cor:extreme} says that $f$ attains a maximum value and a minimum value on $[a,b]$. If both of those values equal $f(a)$, then $f$ is constant, because $f(a)=f(b)$, and the derivative is $0$ everywhere it exists. Otherwise at least one of the two extreme values is attained at a point $c\in\,]a,b[$, and (a) gives $f'(c)=0$.

(c) Apply (b) to
\[g(x)=f(x)-f(a)-\frac{f(b)-f(a)}{b-a}(x-a).\]
Then $g(a)=g(b)=0$, and $g'(x)=f'(x)-\frac{f(b)-f(a)}{b-a}$. A point where $g'$ vanishes is the point $c$ in the statement.

(d) Take $x<y$ in $I$. Part~(c), applied on $[x,y]$, gives $f(y)-f(x)=f'(c)(y-x)$ with $c$ interior to $I$. Each of the three claims is that identity, read with the corresponding hypothesis on $f'$.
\end{proof}

With Theorem~\ref{thm:mvt} in hand, the results of Volume~1 that called on the mean value theorem now rest on a complete proof. That list includes Taylor's formula with remainder, the monotonicity of $\exp$, l'H\^opital's rule, and the inequalities for convex functions. The estimate $|\cos x-\cos y|\le\sin1\,|x-y|$ used in Chapter~\ref{ch:metric} is part~(d), with $M=\sin1$.

\begin{beauty}
Compactness turns a fact that holds at each point into a fact that holds with one constant. You produce a number, or a neighbourhood, at each point. Those neighbourhoods cover the space. A finite subcover leaves finitely many numbers, and the worst of them works everywhere. Heine's theorem is this move for continuity. The positive lower bound in the proof of Corollary~\ref{cor:norms} is the same move for a positive continuous function. A cover by open sets becomes a finite list for the same reason.
\end{beauty}

\section{Tychonoff's theorem and the Cantor space}

Proposition~\ref{prop:finiteproduct} stops at finitely many factors. For an arbitrary product, the finite intersection property and the axiom of choice are the tools that work.

How can a product of uncountably many spaces be compact? The space $[0,1]^{\R}$ is the set of all functions from $\R$ to $[0,1]$. Sequences do not see that topology, and no metric describes it. In 1930, at the age of twenty-four, Andrey Tychonoff proved the theorem for products of segments. The general statement, for an arbitrary product of compact spaces, is from 1935. It needs the axiom of choice, and Kelley's theorem shows that some form of choice is unavoidable. The proof below uses Zorn's lemma. It builds a maximal family of sets that still looks as if the sets have a point in common, and then it produces the point.

\begin{theorem}[Tychonoff, 1935]\label{thm:tychonoff}
The product $\prod_{i\in I}X_i$ of an arbitrary family of compact spaces is compact.
\end{theorem}
\begin{proof}
Write $X=\prod_i X_i$, and let $p_i:X\to X_i$ be the projections. By the closed-set form of compactness at the start of the chapter, it is enough to prove a slightly stronger claim. If $\mathcal A$ is a family of subsets of $X$ with the finite intersection property, some point $x$ lies in the closure of every $A\in\mathcal A$. Applied to a family of closed sets, that point lies in every set of the family.

\emph{A maximal family.} Order by inclusion the families $\mathcal D\supseteq\mathcal A$ of subsets of $X$ that still have the finite intersection property. The union of a chain of such families has the property as well: any finite subcollection from the union already sits inside one family of the chain. Zorn's lemma, Theorem~\ref{thm:zorn}, supplies a maximal family $\mathcal D$. This is where the axiom of choice enters. Maximality gives two closure properties.
\begin{enumerate}[label=(\arabic*)]
\item A finite intersection of members of $\mathcal D$ still belongs to $\mathcal D$. Adding it preserves the finite intersection property, so maximality forces it to be there already.
\item If a set $B$ meets every member of $\mathcal D$, then $B\in\mathcal D$. By (1), the set $B$ meets every finite intersection of members, so adding $B$ preserves the finite intersection property.
\end{enumerate}

\emph{The point.} Fix a coordinate $i$. The sets $p_i(D)$, as $D$ runs through $\mathcal D$, have the finite intersection property in $X_i$. Compactness of $X_i$ puts a point in every one of their closures. Choose one such point and call it $x_i$. Choosing, once for each $i$, is another use of choice. Set $x=(x_i)_{i\in I}$.

\emph{The point lies in every closure.} Take $D\in\mathcal D$, and let $W=\bigcap_{i\in F}p_i^{-1}(U_i)$ be a basic open neighbourhood of $x$. Here $F$ is finite, each $U_i$ is open in $X_i$, and $x_i\in U_i$. For $i\in F$ and any $D'\in\mathcal D$, the set $U_i$ meets $p_i(D')$, because $x_i$ lies in the closure of $p_i(D')$. Thus $p_i^{-1}(U_i)$ meets $D'$. Property~(2) puts $p_i^{-1}(U_i)$ into $\mathcal D$. Property~(1) puts the finite intersection $W$ into $\mathcal D$. The finite intersection property says that $W$ meets $D$. Every basic neighbourhood of $x$ meets $D$, so $x\in\bar D$. Since $\mathcal A\subseteq\mathcal D$, the point $x$ lies in the closure of every member of $\mathcal A$.
\end{proof}

\begin{remark}
Choice cannot be removed from the general statement. Kelley (1950) showed that Tychonoff's theorem implies the axiom of choice. For a countable product of compact metric spaces, choice of that strength is not needed. Such a product is metrizable by Proposition~\ref{prop:product}(c), and a diagonal extraction, as in the proof of Theorem~\ref{thm:compactmetric}, shows that it is sequentially compact. Exercise~\ref{ex:hilbertcube} writes the argument for the Hilbert cube.
\end{remark}

The compact product that this series uses most often is the space of infinite sequences of coin tosses.

\begin{definition}
The \emph{Cantor space} is $\{0,1\}^\N$, where $\N=\{1,2,\dots\}$ indexes the tosses, with the product of the discrete topologies on the factors. The \emph{Cantor set} is
\[C=\Bigl\{\sum_{k\ge1}\frac{2a_k}{3^k}:\ a\in\{0,1\}^\N\Bigr\}\subseteq[0,1],\]
the set of numbers in $[0,1]$ that have a ternary expansion using only the digits $0$ and $2$.
\end{definition}

A basic open set of $\{0,1\}^\N$ is a \emph{cylinder}: the set $[w]$ of all sequences that begin with a fixed finite word $w$. The cylinders are the nodes of the binary tree in Figure~\ref{fig:cantortree}. A point of $\{0,1\}^\N$ is an infinite branch of that tree.

\begin{figure}[ht]
\centering
\fig{ch04-cantortree}
\caption{The Cantor space, drawn as a binary tree of Cantor intervals. A word $w$ of length $n$ names the cylinder $[w]$ and one of the $2^n$ intervals of length $3^{-n}$ in the stage $C_n$. The coral branch through stage $C_4$ reads $0,1,1,0$, and the nested coral intervals shrink to one point of $C$.}
\label{fig:cantortree}
\end{figure}

\begin{proposition}\label{prop:cantor}
\begin{enumerate}[label=(\alph*)]
\item The space $\{0,1\}^\N$ is compact. The formula $\delta(a,b)=2^{-n}$, where $n$ is the first index with $a_n\ne b_n$, and $\delta(a,a)=0$, is a metric for the product topology.
\item The map $h(a)=\sum_k 2a_k 3^{-k}$ is a homeomorphism of $\{0,1\}^\N$ onto $C$. In particular $C$ is compact. Moreover $C=\bigcap_n C_n$, where $C_n$ is the union of the $2^n$ intervals $I_w=[h(w00\ldots),h(w11\ldots)]$, each of length $3^{-n}$, as $w$ runs over the words of length $n$.
\item The map $g(a)=\sum_k a_k 2^{-k}$ is a continuous surjection from $\{0,1\}^\N$ onto $[0,1]$. The composite $g\circ h^{-1}$ is therefore a continuous surjection from $C$ onto $[0,1]$.
\end{enumerate}
\end{proposition}
\begin{proof}
(a) Each factor $\{0,1\}$ is finite, hence compact, so compactness of the product is Tychonoff's theorem. The open $\delta$-ball of radius $2^{-n}$ about $a$ is the cylinder of the first $n$ coordinates of $a$. The $\delta$-open sets are therefore the unions of cylinders, which are exactly the open sets of the product topology. The function $\delta$ is an ultrametric. If $a$ and $b$ agree through index $n$, and $b$ and $c$ agree through index $n$, then $a$ and $c$ agree through index $n$.

(b) If $a$ and $b$ agree through index $n$, then
\[|h(a)-h(b)|\le\sum_{k>n}2\cdot 3^{-k}=3^{-n},\]
so $h$ is continuous. If $n$ is the first index with $a_n\ne b_n$, the same comparison used in Theorem~\ref{thm:continuum} gives
\[|h(a)-h(b)|\ge 2\cdot 3^{-n}-3^{-n}>0.\]
Thus $h$ is injective. Proposition~\ref{prop:compactbasic}(c) and (d) say that $h$ is a homeomorphism onto its image, and that image $C$ is compact. If $a$ begins with the word $w$ of length $n$, then $h(w00\ldots)\le h(a)\le h(w11\ldots)$. The two endpoints differ by $\sum_{k>n}2\cdot 3^{-k}=3^{-n}$, so $C\subseteq C_n$. The intervals $I_{w0}$ and $I_{w1}$ are the left and right thirds of $I_w$. Intervals belonging to distinct words of the same length are therefore disjoint. If $x\in\bigcap_n C_n$, then for each $n$ the point $x$ lies in exactly one interval $I_{w_n}$ with $|w_n|=n$, and $I_{w_{n+1}}\subseteq I_{w_n}$ forces $w_{n+1}$ to extend $w_n$. The words define a sequence $a$. Both $x$ and $h(a)$ lie in $I_{w_n}$, so $|x-h(a)|\le 3^{-n}$ for every $n$. Hence $x=h(a)\in C$.

(c) The same tail estimate, with base $2$ in place of base $3$, shows that $g$ is continuous. Every $x\in[0,1]$ has a binary expansion. Let $a_k$ be the $k$-th digit of the greedy expansion: $a_k=1$ exactly when $\sum_{j<k}a_j 2^{-j}+2^{-k}\le x$. Then $0\le x-\sum_{j\le k}a_j 2^{-j}\le 2^{-k}$, so $g(a)=x$.
\end{proof}

The Cantor set is compact, and it is uncountable, because it is in bijection with $\{0,1\}^\N$. The stage $C_n$ is a union of $2^n$ intervals of length $3^{-n}$, so its total length is $(\frac23)^n$. That length tends to $0$. Chapter~\ref{ch:measures} gives the words ``length zero'' their measure-theoretic meaning. The argument here uses only the lengths of finitely many intervals. The surjection $g\circ h^{-1}$ extends to the Cantor staircase constructed in Chapter~\ref{ch:metric} (Exercise~\ref{ex:devil}). It is also the first step toward a continuous surjection from $[0,1]$ onto $[0,1]^2$. Cantor's bijection in Chapter~\ref{ch:sets} identifies the segment with the square as sets, and it is not continuous. Exercise~\ref{ex:lebesguecurve} builds a continuous surjection.

\section{Arzel\`a--Ascoli}\label{sec:ascoli}

From here on, until the next section, $K$ is a compact metric space and $C(K)$ is the space of continuous real or complex functions on $K$, with the supremum norm. Proposition~\ref{prop:complete}(c) says that $C(K)$ is a Banach space. Continuous functions on $K$ are bounded by Corollary~\ref{cor:extreme}. A subset of a metric space is \emph{relatively compact} when its closure is compact.

\begin{definition}
A family $\mathcal F\subseteq C(K)$ is \emph{equicontinuous} at a point $x\in K$ when, for every $\varepsilon>0$, some $\delta>0$ satisfies $|f(y)-f(x)|<\varepsilon$ for every $f\in\mathcal F$ and every $y$ with $d(x,y)<\delta$. One $\delta$ has to serve the whole family. The family is \emph{equicontinuous} when it is equicontinuous at every point. It is \emph{pointwise bounded} when $\sup_{f\in\mathcal F}|f(x)|<\infty$ for each $x$.
\end{definition}

\begin{figure}[ht]
\centering
\fig{ch04-equicont}
\caption{On the left, $x\mapsto x^n$ on $[0,1]$ fails equicontinuity at $1$: each function is continuous, and the $\delta$ each one needs at $1$ tends to $0$, so no subsequence converges uniformly. On the right, $x\mapsto\frac12\sin(x+c)+\frac c4$ for $0\le c\le 2$ is a $\frac12$-Lipschitz family, bounded as well, so Arzel\`a--Ascoli says the closure is compact.}
\label{fig:equicont}
\end{figure}

Figure~\ref{fig:equicont} is the contrast the theorem measures. Each power $x^n$ is continuous on $[0,1]$. Near $x=1$ the powers become steeper as $n$ grows, and no single $\delta$ keeps all of them inside a fixed $\varepsilon$. The shifted sines on the right share one Lipschitz constant, so one $\delta=\varepsilon/(1/2)$ works for every curve in the family.

\begin{theorem}[Arzel\`a--Ascoli]\label{thm:ascoli}
A family $\mathcal F\subseteq C(K)$ is relatively compact if and only if it is equicontinuous and pointwise bounded.
\end{theorem}
\begin{proof}
($\Leftarrow$) The space $C(K)$ is complete, so the closure $\bar{\mathcal F}$ is complete by Proposition~\ref{prop:complete}(a). Theorem~\ref{thm:compactmetric} says that a complete set is compact once it is totally bounded. An $\varepsilon$-net of $\mathcal F$ is a $2\varepsilon$-net of $\bar{\mathcal F}$, so it is enough to show that $\mathcal F$ itself is totally bounded.

Fix $\varepsilon>0$. At each $x\in K$, equicontinuity supplies a radius $\delta_x$. Finitely many of the balls $B(x_j,\delta_{x_j})$, indexed by $1\le j\le m$, cover $K$. The value vectors
\[v(f)=\bigl(f(x_1),\dots,f(x_m)\bigr),\qquad f\in\mathcal F,\]
lie in a bounded subset of $\R^m$ or $\C^m$, for the norm $\|\cdot\|_\infty$. That bounded set sits in a compact cube, so it is totally bounded. Hence there are functions $f_1,\dots,f_r\in\mathcal F$ such that every $f\in\mathcal F$ has some index $l$ with $\|v(f)-v(f_l)\|_\infty<2\varepsilon$. The factor $2$ is the usual one for subsets: a subset $E$ of a totally bounded space is totally bounded, because if balls of radius $\varepsilon$ about $c_1,\dots,c_p$ cover the ambient space, you pick one point of $E$ in each ball that meets $E$, and the balls of radius $2\varepsilon$ about those points cover $E$. Now take $y\in K$ and choose $j$ with $d(y,x_j)<\delta_{x_j}$. Then
\[\begin{aligned}|f(y)-f_l(y)|&\le|f(y)-f(x_j)|+|f(x_j)-f_l(x_j)|\eqbr+|f_l(x_j)-f_l(y)|<4\varepsilon.\end{aligned}\]
So $\|f-f_l\|_\infty<4\varepsilon$. The functions $f_l$ form a $4\varepsilon$-net of $\mathcal F$.

($\Rightarrow$) If $\bar{\mathcal F}$ is compact, it is bounded, so $\mathcal F$ is pointwise bounded. Fix $x\in K$ and $\varepsilon>0$. Take an $\varepsilon$-net $f_1,\dots,f_r$ of $\mathcal F$, and choose $\delta>0$ so that $|f_l(y)-f_l(x)|<\varepsilon$ whenever $d(x,y)<\delta$, for every one of the finitely many indices $l$. Given $f\in\mathcal F$, choose $l$ with $\|f-f_l\|_\infty<\varepsilon$. Then $|f(y)-f(x)|<3\varepsilon$ whenever $d(x,y)<\delta$.
\end{proof}

\begin{corollary}\label{cor:lipcompact}
Fix $L,M\ge0$ and a point $x_0\in K$. The set of all $L$-Lipschitz functions $f:K\to\R$ with $|f(x_0)|\le M$ is compact in $C(K)$. In particular, every sequence of such functions has a uniformly convergent subsequence.
\end{corollary}
\begin{proof}
The family is equicontinuous: take $\delta=\varepsilon/L$ when $L>0$, and take any $\delta>0$ when $L=0$. A compact metric space is bounded, so the diameter of $K$ is finite, and every such $f$ is bounded by $M+L\operatorname{diam} K$. The family is closed. A uniform limit of $L$-Lipschitz functions is $L$-Lipschitz, and the same is true of a pointwise limit, and the bound $|f(x_0)|\le M$ survives the limit. Theorem~\ref{thm:ascoli} applies.
\end{proof}

The classical use of Arzel\`a--Ascoli is to produce a solution as a limit of approximate solutions, with no uniqueness and no closed formula. Peano's existence theorem for ordinary differential equations is the model, worked in Exercise~\ref{ex:peano}. The same extraction builds Brownian motion as a limit of random walks in Volume~12. Nothing in this chapter depends on that construction.

\section{Stone--Weierstrass}

In 1885 Karl Weierstrass proved that every continuous function on a compact interval $[a,b]$ is a uniform limit of polynomials. The result is surprising once you have seen how wild a continuous function can be. Chapter~\ref{ch:baire} shows that a typical continuous function has a derivative nowhere. Polynomials are the tamest functions on the line. They still come as close as you like, in the uniform distance, to every continuous function. In 1937, and in a simpler form in 1948, Marshall Stone found the reason. Polynomials form an algebra, and they separate points. Many other families do the same. Here $K$ is any compact space, metric or not, and $C(K,\R)$ and $C(K,\C)$ carry the supremum norm. A uniform limit of continuous functions on $K$ is continuous. The proof of Proposition~\ref{prop:complete}(c) gives that fact, and that proof never used a metric on $K$.

\begin{definition}
A set $\mathcal A$ of functions on $K$ \emph{separates points} when, for all $x\ne y$ in $K$, some $f\in\mathcal A$ has $f(x)\ne f(y)$. A \emph{subalgebra} of $C(K,\R)$ is a linear subspace stable under multiplication.
\end{definition}

The proof needs one concrete approximation before the abstract argument. Polynomials have to be able to produce $|t|$ on $[-1,1]$.

\begin{lemma}\label{lem:abs}
Set $p_0=0$ and $p_{n+1}(t)=p_n(t)+\frac12\bigl(t^2-p_n(t)^2\bigr)$. Every $p_n$ is a polynomial, and $0\le|t|-p_n(t)\le\frac2{en}$ holds for all $t\in[-1,1]$ and all $n\ge1$.
\end{lemma}
\begin{proof}
Put $s=|t|$, so $s\in[0,1]$. Expanding the recurrence, and using $t^2=s^2$, gives
\[s-p_{n+1}=(s-p_n)\Bigl(1-\frac{s+p_n}2\Bigr).\]
Suppose $0\le p_n\le s$. The second factor then lies in $[1-s,1]\subseteq[0,1]$, so $0\le s-p_{n+1}\le s-p_n$, and also $p_{n+1}\ge p_n\ge0$. Starting from $p_0=0$, induction gives $0\le p_n\le s$ for every $n$. The same identity, with $p_n\ge0$, yields $s-p_{n+1}\le(s-p_n)\bigl(1-\frac s2\bigr)$, and therefore
\[0\le s-p_n\le s\Bigl(1-\frac s2\Bigr)^n\le se^{-ns/2}\le\frac2{en},\]
where $1-u\le e^{-u}$ was used with $u=s/2$. For $n\ge2$ the function $s\mapsto se^{-ns/2}$ on $[0,1]$ attains its maximum at the critical point $s=\frac2n$, and the value there is $\frac2{en}$. For $n=1$ that critical point is $s=2$, which lies outside $[0,1]$. On $[0,1]$ the function $s\mapsto se^{-s/2}$ is increasing, so its maximum is at $s=1$, and $e^{-1/2}\le\frac2e$.
\end{proof}

\begin{theorem}[Stone--Weierstrass]\label{thm:sw}
\begin{enumerate}[label=(\alph*)]
\item (Real form) Let $\mathcal A\subseteq C(K,\R)$ be a subalgebra that contains the constants and separates points. Then $\mathcal A$ is dense: every $f\in C(K,\R)$ is a uniform limit of elements of $\mathcal A$.
\item (Complex form) Let $\mathcal A\subseteq C(K,\C)$ be a subalgebra that contains the constants, separates points, and is stable under complex conjugation. Then $\mathcal A$ is dense in $C(K,\C)$.
\end{enumerate}
\end{theorem}
\begin{proof}
(a) \emph{Step 1: the closure is a subalgebra, and it is stable under absolute value, max and min.} If $f_n\to f$ and $g_n\to g$ uniformly, then $f_n+g_n\to f+g$ and $cf_n\to cf$. For products,
\[\|f_ng_n-fg\|\le\|f_n\|\|g_n-g\|+\|g\|\|f_n-f\|\to0.\]
The closure is a subalgebra and contains the constants, so it contains every polynomial expression in any one of its members. If $f\in\bar{\mathcal A}$ and $M=\|f\|_\infty>0$, then $Mp_n(f/M)$ lies in $\bar{\mathcal A}$. Lemma~\ref{lem:abs} gives $\bigl|Mp_n(f/M)-|f|\bigr|\le\frac{2M}{en}$. Hence $|f|\in\bar{\mathcal A}$. If $M=0$, then $f=0$ and $|f|$ is already in the algebra. The identities
\[\max(f,g)=\frac12\bigl(f+g+|f-g|\bigr),\qquad\min(f,g)=\frac12\bigl(f+g-|f-g|\bigr)\]
then put $\max(f,g)$ and $\min(f,g)$ into $\bar{\mathcal A}$.

\emph{Step 2: two-point interpolation.} Let $x\ne y$, and let $a,b\in\R$. Some $u\in\mathcal A$ has $u(x)\ne u(y)$, because $\mathcal A$ separates points. The function
\[h=a+(b-a)\frac{u-u(x)}{u(y)-u(x)}\]
lies in $\mathcal A$, since the constants are present, and it satisfies $h(x)=a$ and $h(y)=b$.

\emph{Step 3: a function that stays under $f+\varepsilon$ and touches $f$ at a chosen point.} Let $f\in C(K,\R)$ and $\varepsilon>0$, and fix $x\in K$. For each $y\in K$, Step~2 supplies $h_y\in\mathcal A$ with $h_y(x)=f(x)$ and $h_y(y)=f(y)$. The set $V_y=\{h_y<f+\varepsilon\}$ is open and contains both $y$ and $x$. Finitely many of these sets, $V_{y_1},\dots,V_{y_m}$, cover $K$. The minimum $g_x=\min(h_{y_1},\dots,h_{y_m})$ lies in $\bar{\mathcal A}$, satisfies $g_x<f+\varepsilon$ everywhere, and satisfies $g_x(x)=f(x)$.

\emph{Step 4: push the lower bound up.} The set $W_x=\{g_x>f-\varepsilon\}$ is open and contains $x$. Finitely many of these sets cover $K$. The maximum $g=\max(g_{x_1},\dots,g_{x_r})$ lies in $\bar{\mathcal A}$ and satisfies $f-\varepsilon<g<f+\varepsilon$ at every point. The continuous function $|f-g|$ attains its maximum, and that maximum is still strictly less than $\varepsilon$, so $\|f-g\|_\infty<\varepsilon$. Thus $f\in\bar{\mathcal A}$.

(b) Let $\mathcal A_\R$ be the set of real-valued functions in $\mathcal A$. If $f\in\mathcal A$, then $\operatorname{Re}f=\frac12(f+\bar f)$ and $\operatorname{Im}f=\frac1{2i}(f-\bar f)$ both lie in $\mathcal A_\R$, because $\mathcal A$ is stable under conjugation. So $\mathcal A_\R$ is a real subalgebra, and it contains the real constants. It separates points: if $f(x)\ne f(y)$, then either the real part or the imaginary part separates $x$ from $y$. Part~(a) says that $\mathcal A_\R$ is dense in $C(K,\R)$. Every $f=u+iv\in C(K,\C)$ is then a uniform limit of sums $u_n+iv_n$ with $u_n,v_n\in\mathcal A_\R$.
\end{proof}

\begin{think}
Why does the complex form ask for stability under conjugation? Try to approximate $z\mapsto\bar z$ on the unit circle by polynomials in $z$ alone.
\end{think}

The approximation is impossible. For a polynomial $p(z)=\sum_{k=0}^d c_k z^k$,
\[\int_0^{2\pi}p(e^{it})\,e^{it}\,dt=\sum_kc_k\int_0^{2\pi}e^{i(k+1)t}\,dt=0,\]
because $e^{i(k+1)t}$ has the primitive $\frac{e^{i(k+1)t}}{i(k+1)}$, and that primitive returns to the same value after an interval of length $2\pi$. If polynomials $p_n(z)$ converged to $\bar z$ uniformly on the circle, the integrals against $e^{it}$ would converge to
\[\int_0^{2\pi}e^{-it}e^{it}\,dt=2\pi.\]
They cannot, since each of them is $0$. Polynomials in $z$ do separate the points of the circle, and they contain the constants. The obstruction is conjugation. By the maximum principle in Volume~5, a uniform limit on the circle of polynomials in $z$ extends to a continuous function on the closed disc that is holomorphic inside the disc. That is a much smaller class than $C(S^1,\C)$. Conjugation is the step that puts $\operatorname{Re}f$ and $\operatorname{Im}f$ into the algebra.

\begin{corollary}\label{cor:weierstrass}
\begin{enumerate}[label=(\alph*)]
\item (Weierstrass, 1885) If $K\subseteq\R^n$ is compact, every continuous $f:K\to\R$ is a uniform limit of polynomials in the coordinates $x_1,\dots,x_n$.
\item (Trigonometric form) Every continuous $2\pi$-periodic function $f:\R\to\C$ is a uniform limit of trigonometric polynomials $\sum_{|k|\le N}c_ke^{ikt}$.
\item The space $C([a,b])$ is separable. The polynomials with rational coefficients are dense.
\end{enumerate}
\end{corollary}
\begin{proof}
(a) The polynomials form a subalgebra, they contain the constants, and the coordinate functions separate points of $K$.

(b) By Theorem~\ref{thm:exp}(e), the map $t\mapsto e^{it}$ sends $[0,2\pi]$ continuously onto the circle $S^1$, and it identifies only the endpoints $0$ and $2\pi$. A $2\pi$-periodic $f$ therefore defines a function $F:S^1\to\C$ by $f(t)=F(e^{it})$. This $F$ is continuous. If it failed to be continuous at $e^{is}$, there would be a sequence $e^{it_n}\to e^{is}$ with $|F(e^{it_n})-F(e^{is})|\ge\varepsilon$. The angles $t_n$ may be taken in $[0,2\pi]$. A subsequence, still written $t_n$, converges to some $t\in[0,2\pi]$ with $e^{it}=e^{is}$. Continuity and periodicity of $f$ give $f(t_n)\to f(t)$. The common value of $f$ at angles that represent the same point of the circle is $F(e^{is})$, so $F(e^{it_n})\to F(e^{is})$, which is the contradiction. The polynomials in $z$ and $\bar z$ form a subalgebra of $C(S^1,\C)$. They contain the constants, they are stable under conjugation, and $z$ separates points of the circle. On $S^1$ one has $z^k\bar z^l=e^{i(k-l)t}$. Theorem~\ref{thm:sw}(b) applies.

(c) Given $\varepsilon>0$, part~(a) supplies a polynomial $p$ with $\|f-p\|_\infty<\varepsilon$ on $[a,b]$. Replace the coefficients of $p$ by rationals. Changing the coefficient of $x^k$ by a real number $\eta$ changes the values of $p$ on $[a,b]$ by at most $\eta\max(|a|,|b|)^k$.
\end{proof}

Weierstrass's own argument convolved $f$ with a Gaussian approximate identity. That construction returns in Chapter~\ref{ch:lp}. Bernstein's proof of 1912 is probabilistic and gives the polynomials explicitly. It replaces $f(x)$ by $\mathbb E\bigl[f(\frac{S_n}n)\bigr]$, where $S_n$ counts the heads in $n$ tosses of a coin that lands heads with probability $x$. Exercise~\ref{ex:bernstein} carries out the estimate. Part~(b) of the corollary is the starting point of Fourier analysis in $L^2$, taken up in Volume~8. Fej\'er's theorem, in Chapter~\ref{ch:lp}, gives trigonometric approximants with an explicit formula. The density itself is already proved.

\begin{remark}
Stone--Weierstrass is why a neural network with one hidden layer can approximate continuous functions. On a compact set $K\subseteq\R^n$, the functions $x\mapsto\sum_j c_j\exp(w_j\cdot x)$ form an algebra, because $e^{u}e^{v}=e^{u+v}$. They contain the constants, by taking $w=0$, and they separate points. Theorem~\ref{thm:sw} says they are dense in $C(K,\R)$. Cybenko (1989) and Hornik (1991) proved the same density for sums $\sum_j c_j\sigma(w_j\cdot x+b_j)$, for any sigmoid $\sigma$. For the ReLU $\sigma(t)=\max(t,0)$ in one variable, those sums are exactly the continuous piecewise-linear functions. Heine's theorem makes them dense on a compact interval: uniform continuity supplies a mesh on which the piecewise-linear interpolant stays within $\varepsilon$. Density says nothing about how many terms are needed, and nothing about whether a training procedure finds them.
\end{remark}

\section{The next question}

The cases that worked in the opening are the spaces in which every open cover has a finite subcover. A closed bounded interval has that property, and so does every closed bounded subset of $\R^n$, by Theorem~\ref{thm:heineborel}. An open interval does not. On a nonempty compact space, a continuous real function attains its maximum, by Corollary~\ref{cor:extreme}. On a compact metric space, a continuous map into a metric space is uniformly continuous, by Theorem~\ref{thm:heine}. In a metric space you can test compactness on sequences, or by asking for completeness together with a finite $\varepsilon$-net at every scale. That is Theorem~\ref{thm:compactmetric}. Rolle's theorem and the mean value theorem followed as Theorem~\ref{thm:mvt}, so the arguments of Volume~1 that used them now have a complete foundation.

A bounded sequence in $C([0,1])$ can fail to have a uniformly convergent subsequence, because the closed unit ball is not compact. The tents of height $1$ on the intervals $[\frac1{n+1},\frac1n]$ stay at distance $1$ from each other. Theorem~\ref{thm:ascoli} says what must be added: equicontinuity, together with a bound at each point. The cover property reaches as far as a product can ask. Every product of compact spaces is compact, by Theorem~\ref{thm:tychonoff}, and the Cantor space $\{0,1\}^\N$ of coin tosses is one such product, by Proposition~\ref{prop:cantor}. On a compact space, an algebra of functions that separates points and contains the constants approximates every continuous function, by Theorem~\ref{thm:sw}.

Weierstrass's theorem says that polynomials come within $\varepsilon$ of any continuous function you name. It does not say what a continuous function usually looks like. Take one ``at random''. Does it have a derivative at even a single point? Chapter~\ref{ch:baire} gives that question a precise meaning. The answer is no.

\section{Exercises}

\begin{exercise}[$\star$]\label{ex:dini}
\emph{Dini's theorem.} Let $K$ be compact, and let $f_n,f\in C(K,\R)$ satisfy $f_1\ge f_2\ge\cdots$ and $f_n(x)\to f(x)$ at every $x$. Show that the convergence $f_n\to f$ is uniform. Give examples showing that compactness, monotonicity, and continuity of the limit are each necessary.
\end{exercise}

\begin{exercise}[$\star$]\label{ex:moments}
Let $f:[0,1]\to\R$ be continuous, and suppose $\int_0^1 f(x)x^n\,dx=0$ for every $n\ge0$. Show that $f=0$. In words: the moments of a continuous function determine the function.
\end{exercise}

\begin{exercise}[$\star\star$]\label{ex:isometry}
Let $K$ be a compact metric space, and let $f:K\to K$ satisfy $d(f(x),f(y))=d(x,y)$ for all $x,y$. Show that $f$ is surjective. Find a complete bounded metric space that carries an isometry into itself which misses a point.
\end{exercise}

\begin{exercise}[$\star\star$]\label{ex:edelstein}
\emph{Edelstein's fixed point theorem.} Let $K$ be a compact metric space, and let $f:K\to K$ satisfy $d(f(x),f(y))<d(x,y)$ whenever $x\ne y$. Show that $f$ has a unique fixed point $p$, and that $f^n(x)\to p$ for every $x$. Check that $\sin$, on the interval $[0,1]$, satisfies the strict distance inequality and is not a contraction in the sense of Banach's theorem.
\end{exercise}

\begin{exercise}[$\star\star$]\label{ex:hilbertcube}
\emph{The Hilbert cube.} Let $Q=[0,1]^\N$ with the metric $d(x,y)=\sum_k 2^{-k}|x_k-y_k|$. (a) By a diagonal extraction, and without calling on Tychonoff's theorem, show that $Q$ is compact. (b) Let $(X,d)$ be a compact metric space of diameter at most $1$, and let $(a_k)$ be a dense sequence in $X$. Show that $\iota(x)=(d(x,a_k))_{k\ge1}$ embeds $X$ homeomorphically as a compact subset of $Q$. Every compact metric space therefore sits inside the Hilbert cube.
\end{exercise}

\begin{exercise}[$\star\star$]\label{ex:bernstein}
\emph{Bernstein's proof of Weierstrass.} For $f\in C([0,1])$ define $B_nf(x)=\sum_{k=0}^n f(\frac kn)\binom nkx^k(1-x)^{n-k}$. (a) Show that $\sum_k\binom nkx^k(1-x)^{n-k}\bigl(\frac kn-x\bigr)^2=\frac{x(1-x)}n$. (b) Deduce the uniform convergence $B_nf\to f$ on $[0,1]$. (c) Read $B_nf(x)$ as an average over $n$ tosses of a coin that lands heads with probability $x$.
\end{exercise}

\begin{exercise}[$\star\star$]\label{ex:lebesguecurve}
\emph{Lebesgue's space-filling curve.} (a) From the maps in Proposition~\ref{prop:cantor}, build a continuous surjection $\psi:C\to[0,1]^2$. Split each sequence into its odd terms and its even terms. (b) Extend $\psi$ to a map $\Psi:[0,1]\to[0,1]^2$ that is affine on each connected component of $[0,1]\setminus C$, and show that $\Psi$ is continuous. Thus $[0,1]$ maps continuously onto the square. The map is not bijective; that is Exercise~\ref{ex:notR2}.
\end{exercise}

\begin{exercise}[$\star\star\star$]\label{ex:peano}
\emph{Peano's existence theorem.} Let $f:[0,1]\times\R\to\R$ be continuous and bounded by a constant $M$, and let $y_0\in\R$. (a) For each $n\ge1$, let $y_n$ be the Euler polygon: $y_n(0)=y_0$, and on each interval $[\frac kn,\frac{k+1}n]$ the function $y_n$ is affine with slope $f\bigl(\frac kn,y_n(\frac kn)\bigr)$. Show that some subsequence converges uniformly to a limit $y$. (b) Show that $y(t)=y_0+\int_0^t f(s,y(s))\,ds$, so $y$ is a $C^1$ solution of $y'=f(t,y)$ with $y(0)=y_0$. (c) Take $f(t,y)=\sqrt{\min(|y|,1)}$ and $y_0=0$, and exhibit two solutions. Existence does not need the Lipschitz condition of Picard--Lindel\"of (Exercise~\ref{ex:picard}). Uniqueness does.
\end{exercise}

\begin{exercise}[$\star$]\label{ex:v7c04a}
Let $A$ and $B$ be compact subsets of a topological space $X$. Show that $A\cup B$ is compact. The space $X$ is not assumed Hausdorff.
\end{exercise}

\begin{exercise}[$\star\star$]\label{ex:v7c04b}
Let $X=(0,1]$ with the subspace metric inherited from $\R$, and let $x_n=\frac1n$. Show that no subsequence of $(x_n)$ converges in $X$. Conclude that $X$ is not compact, and that the failure is incompleteness rather than a failure of total boundedness.
\end{exercise}

\begin{exercise}[$\star\star$]\label{ex:v7c04c}
Let $K$ be a nonempty compact space, and let $f:K\to\R$ be continuous. Write $m=\sup_K f$ and $M=\{x\in K:f(x)=m\}$. Show that $M$ is nonempty and compact.
\end{exercise}

\begin{exercise}[$\star\star\star$]\label{ex:v7c04d}
Let $\mathcal H$ be the set of functions $f:[0,1]\to\R$ such that $|f(x)|\le 1$ for every $x$, and $|f(x)-f(y)|\le\sqrt{|x-y|}$ for all $x,y$. Show that $\mathcal H$ is a compact subset of $C([0,1])$.
\end{exercise}

\section{Hints}

\begin{hint}{ex:dini}
For $\varepsilon>0$, look at the closed sets $F_n=\{f_n-f\ge\varepsilon\}$. They decrease, and their total intersection is empty. For the examples, compare $x^n$ on $[0,1[$ with $x^n$ on $[0,1]$, and remember the tents from the unit ball.
\end{hint}
\begin{hint}{ex:moments}
Linearity gives $\int_0^1 fp=0$ for every polynomial $p$. Approximate $f$ itself by polynomials, and look at $\int_0^1 f^2$.
\end{hint}
\begin{hint}{ex:isometry}
The image $f(K)$ is compact, hence closed. If $y$ misses it, the distance $\varepsilon$ from $y$ to $f(K)$ is positive. Watch the orbit $y,f(y),f^2(y),\dots$.
\end{hint}
\begin{hint}{ex:edelstein}
The function $x\mapsto d(x,f(x))$ is continuous on a compact space, so it attains a minimum. For the orbits, $d(f^n(x),p)$ is nonincreasing. Look at a limit point of the orbit.
\end{hint}
\begin{hint}{ex:hilbertcube}
(a) Pass to a subsequence along which the first coordinates converge, then to a further subsequence along which the second coordinates converge, and continue. The diagonal subsequence converges in every coordinate, and the weights $2^{-k}$ make the tail small. (b) Each coordinate of $\iota$ is $1$-Lipschitz. To separate $x\ne y$, choose $a_k$ very close to $x$.
\end{hint}
\begin{hint}{ex:bernstein}
(a) Differentiate $(x+y)^n=\sum\binom nkx^ky^{n-k}$ in $x$, twice, or compute $\sum k\binom nk$ and $\sum k(k-1)\binom nk$ directly. (b) Split the sum into the terms with $|\frac kn-x|<\delta$ and the rest, and call on Heine for the first group.
\end{hint}
\begin{hint}{ex:lebesguecurve}
(a) The map is $\psi=(g\circ\mathrm{odd},\,g\circ\mathrm{even})\circ h^{-1}$. (b) Only finitely many gaps are longer than a given $\eta>0$. If $x_n\to x\in C$ and the points $x_n$ lie in gaps, the endpoints of those gaps also tend to $x$, unless one single gap is visited infinitely often.
\end{hint}
\begin{hint}{ex:peano}
(a) Each polygon $y_n$ is $M$-Lipschitz and starts at $y_0$. Use Corollary~\ref{cor:lipcompact}. (b) Write $y_n(t)=y_0+\int_0^t f\bigl(\frac{\lfloor ns\rfloor}n,\,y_n(\frac{\lfloor ns\rfloor}n)\bigr)\,ds$, and use the uniform continuity of $f$ on a compact rectangle. (c) Try $y=0$ and $y=\frac{t^2}4$.
\end{hint}
\begin{hint}{ex:v7c04a}
An open cover of the union covers $A$ and covers $B$. Take a finite subcover of each set, and put the two finite lists together.
\end{hint}
\begin{hint}{ex:v7c04b}
Every subsequence still tends to $0$ in $\R$. A limit in $X$ would be a limit in $\R$. For total boundedness, space the centres about $\varepsilon$ apart.
\end{hint}
\begin{hint}{ex:v7c04c}
Corollary~\ref{cor:extreme} puts at least one point into $M$. The set $M$ is the preimage of the closed singleton $\{m\}$.
\end{hint}
\begin{hint}{ex:v7c04d}
The square-root modulus gives one $\delta$ for the whole family: try $\delta=\varepsilon^2$. A uniform limit keeps both the bound $1$ and the square-root inequality. Then use Theorem~\ref{thm:ascoli}.
\end{hint}

\section{Solutions}

\begin{solution}{ex:dini}
\idea Pointwise convergence says that each $x$ is eventually good. Monotonicity says that once $x$ is good, it stays good. Compactness turns ``each point eventually'' into ``all points from one index on'', through the finite intersection property.
\writeup Fix $\varepsilon>0$ and set $F_n=\{x:f_n(x)-f(x)\ge\varepsilon\}$. Each $F_n$ is closed, because $f_n-f$ is continuous. The inequality $f_n\ge f_{n+1}\ge f$ makes the sets decrease: $F_{n+1}\subseteq F_n$. Their intersection is empty, because $f_n(x)\to f(x)$ at every $x$. Compactness of $K$ says that some finite intersection is empty. The sets decrease, so that finite intersection is just some $F_N$, and $F_n=\varnothing$ for all $n\ge N$. Thus $0\le f_n-f<\varepsilon$ everywhere for $n\ge N$.

Three examples show the three hypotheses at work. On $[0,1[$, the functions $x^n$ decrease pointwise to $0$, and the convergence is not uniform, because the domain is not compact. On $[0,1]$, the same functions decrease pointwise to the function that is $0$ on $[0,1[$ and $1$ at $1$, which is discontinuous, so the continuity of the limit was necessary. The tents of height $1$ supported on $[\frac1{n+1},\frac1n]$ tend pointwise to $0$ on the compact interval $[0,1]$, they are continuous, and they are not monotone. The convergence is not uniform, because each tent has supremum norm $1$.
\end{solution}

\begin{solution}{ex:moments}
\idea If $f$ is orthogonal to every polynomial, it is orthogonal to every uniform limit of polynomials. One of those limits is $f$ itself, by Corollary~\ref{cor:weierstrass}.
\writeup Linearity gives $\int_0^1 fp=0$ for every polynomial $p$. Corollary~\ref{cor:weierstrass} supplies polynomials $p_n$ with $p_n\to f$ uniformly on $[0,1]$. Then
\[\Bigl|\int_0^1f^2-\int_0^1fp_n\Bigr|\le\|f\|_\infty\|f-p_n\|_\infty\to0,\]
so $\int_0^1 f^2=0$. A continuous function that is nonnegative and has integral $0$ must vanish. If $f^2(x_0)>0$, continuity keeps $f^2>\frac12 f^2(x_0)$ on a whole interval about $x_0$, and the integral over that interval is positive. Hence $f^2=0$, and $f=0$.
\end{solution}

\begin{solution}{ex:isometry}
\idea An isometry that misses a point pushes that point along an orbit whose terms stay at least $\varepsilon$ apart. A compact space has no infinite $\varepsilon$-separated set.
\writeup The map $f$ is continuous, because it preserves distances. Proposition~\ref{prop:compactbasic}(c) says that $f(K)$ is compact, and part~(b) of the same proposition says that $f(K)$ is closed in $K$. Suppose some $y\in K$ lies outside $f(K)$. The distance $\varepsilon=d(y,f(K))$ is positive, because $f(K)$ is closed. Write $y_n=f^n(y)$. For $m>n$, the iterate $f^n$ is still an isometry, so
\[d(y_n,y_m)=d\bigl(y,f^{m-n}(y)\bigr)\ge\varepsilon,\]
since $f^{m-n}(y)\in f(K)$. No subsequence of $(y_n)$ is Cauchy. That contradicts sequential compactness, Theorem~\ref{thm:compactmetric}.

Compactness was used. On $\N$ with the discrete metric, where distinct points are at distance $1$, the space is complete and bounded. The shift $n\mapsto n+1$ is an isometry, and it misses $0$.
\end{solution}

\begin{solution}{ex:edelstein}
\idea A strict decrease of distances is not a Banach contraction: the ratio may creep up to $1$, and there is no single $k<1$ to iterate. Compactness replaces the missing uniform constant. A continuous function attains its minimum, and an orbit in a compact space has a limit point.
\writeup The function $\phi(x)=d(x,f(x))$ is continuous, because $f$ is $1$-Lipschitz. Corollary~\ref{cor:extreme} says that $\phi$ attains its minimum at some point $p$. If $f(p)\ne p$, then
\[\phi(f(p))=d\bigl(f(p),f(f(p))\bigr)<d(p,f(p))=\phi(p),\]
which contradicts the minimum. So $f(p)=p$. If $q$ were another fixed point, then $d(p,q)=d(f(p),f(q))<d(p,q)$, which is absurd. Thus the fixed point is unique.

Fix $x\in K$, write $x_n=f^n(x)$, and set $r_n=d(x_n,p)$. The hypothesis gives $d(f(x_n),f(p))\le d(x_n,p)$, with a strict inequality unless $x_n=p$. So $(r_n)$ is nonincreasing, and it converges to some $r\ge0$. Sequential compactness supplies a subsequence $x_{n_k}\to q$, and then $d(q,p)=r$. Also $x_{n_k+1}=f(x_{n_k})\to f(q)$, so $d(f(q),p)=\lim r_{n_k+1}=r$. If $r>0$, then $q\ne p$, and $d(f(q),p)=d(f(q),f(p))<d(q,p)=r$, a contradiction. Hence $r=0$, and the whole sequence $x_n$ converges to $p$, because $r_n$ decreases to $0$.

On $[0,1]$, the mean value theorem gives $|\sin x-\sin y|=|\cos c|\,|x-y|$ for some $c$ strictly between $x$ and $y$. For $x\ne y$ that $c$ lies in $]0,1]$, where $0<\cos c<1$. So $\sin$ strictly decreases distances, and $\sin$ maps $[0,1]$ into itself. It is not a Banach contraction. The difference quotient $\frac{\sin x}x$ tends to $1$ as $x\to0$, so no constant $k<1$ works for every pair. The fixed point is $0$.
\end{solution}

\begin{solution}{ex:hilbertcube}
\idea Convergence in $Q$ is coordinatewise convergence. The weights $2^{-k}$ make the tail of the series smaller than any $\varepsilon$ you name, so a diagonal subsequence that settles in each coordinate settles in $Q$. For (b), the distances to a dense sequence are coordinates that know where a point sits.
\writeup (a) Suppose $x^{(n)}_k\to x_k$ for every $k$. Then, for any fixed $K$,
\[d\bigl(x^{(n)},x\bigr)\le\sum_{k\le K}\bigl|x^{(n)}_k-x_k\bigr|+2^{-K}.\]
The sum is $<2^{-K}$ for all large $n$, so the distance is $<2^{1-K}$. Coordinatewise convergence therefore implies convergence in $Q$. Now take a sequence $(x^{(n)})$ in $Q$. Theorem~\ref{thm:heineborel} gives an infinite set $N_1$ of indices along which the first coordinates converge, then an infinite set $N_2\subseteq N_1$ along which the second coordinates converge, and so on. Choose $n_1<n_2<\cdots$ with $n_j\in N_j$. For each fixed $k$, the indices $n_j$ with $j\ge k$ lie in $N_k$, so the $k$-th coordinates of $x^{(n_j)}$ converge. The diagonal subsequence converges in $Q$. Thus $Q$ is sequentially compact, and Theorem~\ref{thm:compactmetric} says it is compact.

(b) The inequality $|d(x,a_k)-d(y,a_k)|\le d(x,y)$ gives $d(\iota(x),\iota(y))\le\sum_k 2^{-k}d(x,y)\le d(x,y)$. So $\iota$ is continuous. The diameter hypothesis puts every coordinate in $[0,1]$, so the image lies in $Q$. If $x\ne y$, let $r=d(x,y)$ and choose $a_k$ with $d(x,a_k)<\frac r3$. Then $d(y,a_k)>\frac{2r}3$, so $\iota(x)\ne\iota(y)$. A continuous injection from a compact space into a Hausdorff space is a homeomorphism onto its image, by Proposition~\ref{prop:compactbasic}(d), and the image is compact. A dense sequence $(a_k)$ exists because a compact metric space is separable.
\end{solution}

\begin{solution}{ex:bernstein}
\idea Under the binomial weights, the ratio $\frac kn$ concentrates near $x$. Its variance is $\frac{x(1-x)}n$, which is at most $\frac1{4n}$. The weighted average of the numbers $f(\frac kn)$ therefore sees only values of $f$ near $x$, and uniform continuity makes those values close to $f(x)$.
\writeup (a) Write $w_k=\binom nkx^k(1-x)^{n-k}$. The binomial theorem gives $\sum_k w_k=1$. The identities $k\binom nk=n\binom{n-1}{k-1}$ and $k(k-1)\binom nk=n(n-1)\binom{n-2}{k-2}$ give $\sum_k kw_k=nx$ and $\sum_k k(k-1)w_k=n(n-1)x^2$. Therefore $\sum_k k^2 w_k=n(n-1)x^2+nx$, and
\[\begin{aligned}\sum_kw_k\Bigl(\frac kn-x\Bigr)^2&=\frac{n(n-1)x^2+nx}{n^2}-2x^2+x^2\eqbr=\frac{x-x^2}n.\end{aligned}\]

(b) Let $\varepsilon>0$. Heine's theorem supplies $\delta>0$ such that $|f(s)-f(x)|<\varepsilon$ whenever $|s-x|<\delta$. Split the average:
\[\begin{aligned}|B_nf(x)-f(x)|&\le\sum_kw_k\Bigl|f\Bigl(\frac kn\Bigr)-f(x)\Bigr|\\&\le\varepsilon+2\|f\|_\infty\sum_{|\frac kn-x|\ge\delta}w_k\eqbr\le\varepsilon+\frac{2\|f\|_\infty}{\delta^2}\cdot\frac1{4n},\end{aligned}\]
where the last step uses $1\le\bigl(\frac kn-x\bigr)^2/\delta^2$ on the second range, together with $x(1-x)\le\frac14$. For $n\ge\|f\|_\infty/(2\delta^2\varepsilon)$, one has $|B_nf-f|\le 2\varepsilon$ everywhere on $[0,1]$.

(c) If $S_n$ is the number of heads in $n$ independent tosses with heads-probability $x$, then $\mathbb P(S_n=k)=w_k$, so $B_nf(x)=\mathbb E\bigl[f(\frac{S_n}n)\bigr]$. Part~(a) says $\operatorname{Var}(\frac{S_n}n)=\frac{x(1-x)}n$. Part~(b) is Chebyshev's inequality, applied uniformly in the bias $x$. The weak law of large numbers, uniform in $x$, is a proof of Weierstrass's theorem.
\end{solution}

\begin{solution}{ex:lebesguecurve}
\idea A point of the Cantor set is an infinite sequence of coin tosses. The odd tosses name one point of $[0,1]$, and the even tosses name another. On a gap of the Cantor set there is nothing to encode, so join the values at the two endpoints by a straight segment.
\writeup (a) The splitting map $s:\{0,1\}^\N\to\{0,1\}^\N\times\{0,1\}^\N$ given by
\[a\mapsto\bigl((a_1,a_3,a_5,\ldots),\,(a_2,a_4,a_6,\ldots)\bigr)\]
is a bijection. It is continuous, because each output coordinate is one input coordinate. Let $g$ and $h$ be the maps of Proposition~\ref{prop:cantor}, and set $\psi=(g\times g)\circ s\circ h^{-1}$. The inverse $h^{-1}$ is continuous by part~(b) of that proposition. The product $g\times g$ is continuous and onto $[0,1]^2$, and $s$ is onto. So $\psi:C\to[0,1]^2$ is a continuous surjection.

(b) The complement $[0,1]\setminus C$ is a countable union of disjoint open intervals $]\alpha,\beta[$, the gaps, with endpoints $\alpha,\beta\in C$. Indeed $[0,1]\setminus C=\bigcup_n\bigl([0,1]\setminus C_n\bigr)$ by Proposition~\ref{prop:cantor}(b), and each interval removed at stage $n$ has its endpoints in $C$. Define $\Psi=\psi$ on $C$, and on a gap $]\alpha,\beta[$ set
\[\Psi(x)=\psi(\alpha)+\frac{x-\alpha}{\beta-\alpha}\bigl(\psi(\beta)-\psi(\alpha)\bigr).\]
This extension is continuous on each open gap. To check continuity at a point $x\in C$, take $x_n\to x$. It is enough to show that every subsequence of $\Psi(x_n)$ has a further subsequence converging to $\Psi(x)$. Every subsequence has a further subsequence of one of the following three kinds.

If all the terms lie in $C$, continuity of $\psi$ gives the limit. If all the terms lie in one fixed gap $]\alpha,\beta[$, then $x$ is one of the endpoints $\alpha,\beta$, because $x\in C$ and $x\in[\alpha,\beta]$. The affine formula extends continuously up to those endpoints. Otherwise, after passing to a subsequence, the points $x_n$ lie in gaps $]\alpha_n,\beta_n[$, each visited only finitely often. Only finitely many gaps have length $\ge\eta$, because stage $n$ removes gaps of length $3^{-n}$. Thus $\beta_n-\alpha_n\to0$. Since $x$ is not inside $]\alpha_n,\beta_n[$, both $|\alpha_n-x|$ and $|\beta_n-x|$ are at most $|x_n-x|+(\beta_n-\alpha_n)$, which tends to $0$. Continuity of $\psi$ sends $\psi(\alpha_n)$ and $\psi(\beta_n)$ to $\psi(x)$. The value $\Psi(x_n)$ lies on the segment joining those two image points, so it tends to $\psi(x)$ as well. Hence $\Psi$ is continuous. It is onto $[0,1]^2$ because $\psi$ already is.
\end{solution}

\begin{solution}{ex:peano}
\idea Euler's polygons are approximate solutions, and they are all $M$-Lipschitz. Arzel\`a--Ascoli extracts a uniformly convergent subsequence. Uniform continuity of $f$ on a compact rectangle lets the limit pass inside the integral. Different subsequences need not share the same limit, and uniqueness can fail.
\writeup (a) Each $y_n$ is continuous and piecewise affine, with slopes bounded by $M$, so $y_n$ is $M$-Lipschitz, and $y_n(0)=y_0$. Corollary~\ref{cor:lipcompact}, applied on $K=[0,1]$, gives a subsequence $y_{n_j}$ converging uniformly to some $y$. The limit is itself $M$-Lipschitz.

(b) On the half-open piece $]\frac kn,\frac{k+1}n[$, the slope of $y_n$ equals $f(\tau_n(s),y_n(\tau_n(s)))$, where $\tau_n(s)=\frac{\lfloor ns\rfloor}n$. Integrating the slope gives
\[y_n(t)=y_0+\int_0^t f\bigl(\tau_n(s),y_n(\tau_n(s))\bigr)\,ds.\]
The integral of a step function is the sum of its pieces. Every value $y_n(s)$, and the limit $y(s)$, lies in the interval $[y_0-M,y_0+M]$. The rectangle $R=[0,1]\times[y_0-M,y_0+M]$ is compact, so Heine's theorem says that $f$ is uniformly continuous on $R$. Also $|\tau_n(s)-s|\le\frac1n$, and
\[\begin{aligned}|y_n(\tau_n(s))-y(s)|&\le|y_n(\tau_n(s))-y_n(s)|+|y_n(s)-y(s)|\eqbr\le\tfrac Mn+\|y_n-y\|_\infty.\end{aligned}\]
Along the subsequence, the integrand converges to $f(s,y(s))$ uniformly in $s$. Passing to the limit,
\[y(t)=y_0+\int_0^t f(s,y(s))\,ds.\]
The integrand is continuous. The fundamental theorem of calculus from Volume~1 says that $y$ is $C^1$, with $y'(t)=f(t,y(t))$ and $y(0)=y_0$.

(c) The zero function is a solution. For $y=\frac{t^2}4$, one has $0\le y\le\frac14$ on $[0,1]$, so $f(t,y)=\sqrt y$, and $y'=\frac t2=\sqrt y$. Both solutions start at $0$. Near $y=0$ the function $f$ is not Lipschitz, because $\sqrt{|y|}/|y|\to\infty$ as $y\to0$.
\end{solution}

\begin{solution}{ex:v7c04a}
\idea Compactness is a statement about covers. A cover of the union is a cover of each piece, and two finite lists make one finite list.
\writeup Let $\mathcal U$ be a cover of $A\cup B$ by open sets of $X$. The same family covers $A$. Compactness of $A$ gives a finite subcollection that covers $A$. The same family covers $B$, so a second finite subcollection covers $B$. The union of those two finite lists is a finite subcover of $A\cup B$. Hausdorff separation was never used.
\end{solution}

\begin{solution}{ex:v7c04b}
\idea The picture is Figure~\ref{fig:rwescape}. The terms run toward a point that $X$ does not contain, and every subsequence does the same. A finite grid still nets the set, so the missing ingredient is the limit, not the net.
\writeup Let $n_k$ be strictly increasing. Then $x_{n_k}=\frac1{n_k}\to 0$ in $\R$. Suppose the subsequence converged to a point $L\in X$. Convergence in the subspace metric is convergence in $\R$ to a point of $X$, so the subsequence would also converge to $L$ in $\R$. Limits in $\R$ are unique, hence $L=0$. But $0\notin X$. No subsequence converges in $X$. Theorem~\ref{thm:compactmetric} says that $X$ is not compact.

The sequence is Cauchy in $X$, and it has no limit in $X$, so $X$ is not complete. It is totally bounded. If $\varepsilon\ge 1$, the single point $1$ is an $\varepsilon$-net of $(0,1]$. If $0<\varepsilon<1$, let $N$ be the greatest integer with $N\varepsilon\le 1$. The points $\varepsilon,2\varepsilon,\dots,N\varepsilon$ lie in $X$. A point of $(0,1]$ smaller than $\varepsilon$ lies within $\varepsilon$ of $\varepsilon$. A point $x$ with $k\varepsilon\le x<(k+1)\varepsilon$ lies within $\varepsilon$ of $k\varepsilon$, and $1$ lies within $\varepsilon$ of $N\varepsilon$ because $1-N\varepsilon<\varepsilon$. Compactness fails because a Cauchy sequence has no limit, not because some scale refuses a finite net.
\end{solution}

\begin{solution}{ex:v7c04c}
\idea The extreme value corollary puts the supremum into the image. The set where that value is attained is a closed subset of $K$, cut out by one equation.
\writeup Corollary~\ref{cor:extreme} says that $f$ attains its maximum, so $M$ is nonempty, and $m$ is a value of $f$ rather than a value the function only approaches. The singleton $\{m\}$ is closed in $\R$. Continuity of $f$ says that $M=f^{-1}(\{m\})$ is closed in $K$. Proposition~\ref{prop:compactbasic}(a) says that a closed subset of a compact space is compact. So $M$ is compact.
\end{solution}

\begin{solution}{ex:v7c04d}
\idea One modulus of continuity for the whole family is equicontinuity. The bound $|f|\le 1$ is pointwise boundedness. The two inequalities survive uniform limits, so the set is closed, and Arzel\`a--Ascoli makes the closure compact.
\writeup Equicontinuity: given $\varepsilon>0$, set $\delta=\varepsilon^2$. If $|x-y|<\delta$, then for every $f\in\mathcal H$,
\[|f(x)-f(y)|\le\sqrt{|x-y|}<\varepsilon.\]
The same $\delta$ serves every function in the family. Pointwise, $|f(x)|\le 1$. Theorem~\ref{thm:ascoli} says that $\mathcal H$ is relatively compact in $C([0,1])$.

The set is closed. Let $f_n\in\mathcal H$ and $\|f_n-f\|_\infty\to 0$. Then $|f(x)|=\lim_n|f_n(x)|\le 1$. For each pair $x,y$,
\[|f(x)-f(y)|=\lim_n|f_n(x)-f_n(y)|\le\sqrt{|x-y|}.\]
The limit $f$ is continuous, because a uniform limit of continuous functions is continuous, as in the proof of Proposition~\ref{prop:complete}(c). Thus $f\in\mathcal H$. A closed relatively compact set is compact, so $\mathcal H$ is compact in $C([0,1])$.
\end{solution}

\section*{Sources and further reading}
\addcontentsline{toc}{section}{Sources and further reading}

\begin{itemize}
\item Compactness, the tube lemma, and the maximal-family proof of Tychonoff's theorem are in J.~Munkres, \emph{Topology}, 2nd ed., Prentice Hall, 2000, Chapters~3 and~5. The argument above follows that pattern.
\item Compact metric spaces, equicontinuity and Stone--Weierstrass are treated in W.~Rudin, \emph{Principles of Mathematical Analysis}, 3rd ed., McGraw--Hill, 1976, Chapters~2 and~7. The iteration in Lemma~\ref{lem:abs} is his Exercise~7.23.
\item Products of segments are treated in A.~Tychonoff, \emph{\"Uber die topologische Erweiterung von R\"aumen}, Mathematische Annalen 102 (1930), 544--561. That the general product theorem yields the axiom of choice is J.~L.~Kelley, \emph{The Tychonoff product theorem implies the axiom of choice}, Fundamenta Mathematicae 37 (1950), 75--76.
\item The polynomial theorem on an interval is K.~Weierstrass, \emph{\"Uber die analytische Darstellbarkeit sogenannter willk\"urlicher Functionen einer reellen Ver\"anderlichen}, Sitzungsberichte der Akademie zu Berlin (1885), 633--639 and 789--805. The algebra form used above is M.~H.~Stone, \emph{The generalized Weierstrass approximation theorem}, Mathematics Magazine 21 (1948), 167--184 and 237--254.
\item The probabilistic proof written out in Exercise~\ref{ex:bernstein} is S.~Bernstein, \emph{D\'emonstration du th\'eor\`eme de Weierstrass fond\'ee sur le calcul des probabilit\'es}, Communications de la Soci\'et\'e math\'ematique de Kharkow 13 (1912), 1--2.
\item Lebesgue's curve, together with the curves of Peano and of Hilbert, is the subject of H.~Sagan, \emph{Space-Filling Curves}, Springer, 1994.
\item The density theorems for sigmoidal sums, quoted after Corollary~\ref{cor:weierstrass}, are G.~Cybenko, \emph{Approximation by superpositions of a sigmoidal function}, Mathematics of Control, Signals and Systems 2 (1989), 303--314, and K.~Hornik, \emph{Approximation capabilities of multilayer feedforward networks}, Neural Networks 4 (1991), 251--257.
\item \emph{For the dates in the opening.} M.~Raman-Sundstr\"om, ``A pedagogical history of compactness'', \emph{American Math. Monthly} 122 (2015), 619--635, traces the definitions through Heine, Borel, Fr\'echet and Alexandroff--Urysohn. The unit-ball theorem of F.~Riesz is in \emph{\"Uber lineare Funktionalgleichungen}, Acta Mathematica 41
(1918), 71--98. Stone's first version is \emph{Applications of the theory of Boolean rings to general topology}, Transactions of the American
Mathematical Society 41 (1937), 375--481. The general product theorem is A.~Tychonoff, \emph{\"Uber einen Funktionenraum}, Mathematische Annalen 111
(1935), 762--766.
\end{itemize}
